\section{Day 4: Spectrum Properties (Sep.\ 17, 2026)}
Recall, from last time, that we took
\[ \abs{\Spec R} = \{p \in R \mid p \text{ is prime}\}, \]
and gave it a topology with a basis of closed sets $Z(f) = \{p \in \Spec R \mid f(p) = 0\}$, where $f \in R$, and $f(p) = 0$ meant the image of $f$ in $k(p) = \Frac(R/p)$. Equivalently, $f \in p$.

\begin{theorem} \label{thm:vanishing_locus_of_ideal}
    If $I = (A)$ is the ideal generated by $A \subset R$, then $z(A) = z(I) = \{p \in \Spec R \mid I \subset p\}$.
\end{theorem}
\begin{proof}
    $z(A) = \{p \in \Spec R \mid A \subset p\} = \{p \in \Spec R \mid I \subset p\} = z(I)$.
\end{proof}

\begin{theorem} \label{thm:ring_to_spectrum_contravariant}
    The assignment $R \to \abs{\Spec R}$ is a contravariant functor $\mathbf{Ring} \to \mathbf{Top}$.
\end{theorem}
We introduce some category theory here. A \textit{category} is a set of objects $\Ob(C)$, a collection of morphisms $\{\Mor(X, Y) : X, Y \in \Ob(C)\}$, and a composition map
\[ \circ \colon \Mor(X, Y) \times \Mor(Y, Z) \to \Mor(X, Z) \]
satisfying the properties,
\begin{enumerate}[(i)]
    \item Composition is associative, i.e., $f \circ (g \circ h) = (f \circ g) \circ h$ whenever this makes sense,
    \item An identity map exists, i.e., $\id_X \in \Mor(X, X)$ such that $\id_X \circ f = f = f \circ \id_X$ whenever $f \in \Mor(X, Y)$.
\end{enumerate}

In the theorem above, there are two categories: rings and topological spaces. $\mathbf{Ring}$ is a category whose objects are rings and morphisms are ring maps (there's nothing to check here because the morphisms are determined at the level of elements). In particular, $\mathbf{Ring}$ is a subcategory of $\mathbf{Set}$. $\mathbf{Top}$ is a category whose objects are topological spaces and morphisms are continuous maps.

A \textit{functor} $F \colon \SCC \to \SCD$ between categories is the data of two maps (both denoted $F$)
\[ F \colon \Ob(\SCC) \to \Ob(\SCD), \quad F \colon \Mor(X, Y) \to \Mor(F(X), F(Y)), \]
satisfying the following relations,
\begin{enumerate}[(i)]
    \item $F(f) \circ F(g) = F(f \circ g)$ when this makes sense,
    \item $F(\id_X) = \id_{F(X)}$.
\end{enumerate}

We now define the notion of an \textit{opposite category}. Let $\SCD$ be a category. Then $\SCD^{\opname{op}}$ as a category consists has $\Ob(\SCD^{\opname{op}}) = \Ob(\SCD)$ and
\[ \Mor_{\SCD^{\opname{op}}}(X, Y) \coloneq \{f^{\opname{op}} \mid f \in \Mor_\SCD(Y, X)\}. \]
Composition works by $f^{\opname{op}} \circ g^{\opname{op}} \coloneq (g \circ f)^{\opname{op}}$. With this, we may now define contravariant functors. A \textit{contravariant functor} is a functor $F \colon \SCC \to \SCD^{\opname{op}}$ with
\[ F \colon \Ob(\SCC) \to \Ob(\SCD), \quad F \colon \Mor(X, Y) \to \Mor(F(Y), F(X)), \]
satisfying obvious rules $F(\id_X) = \id_{F(X)}$ and $F(f \circ g) = F(g) \circ F(f)$. As an example of more functors, we have the forgetful functor $\mathbf{Ring} \to \mathbf{Set}$, which is basically a functor that ``forgets the ring structure''. We also have the contravariant functor from $\mathbf{Vect}_K \to \mathbf{Vect}_K$, where we sent $V \to V^\ast = \Hom(V, K)$, given by the dual map $(f \colon V \to W) \mapsto (f^\ast \colon W^\ast \to V^\ast)$.

\newpage
As another example, consider the functor $\mathbf{Mfld} \to \mathbf{Ring}$, given by $F \colon M \to C^0(M, \RR)$, and $(\Phi \colon M \to N) \mapsto (\Phi^\ast \colon C^0(N, \RR) \to C^0(M, \RR))$, i.e., the pullback.

We now prove \ref{thm:ring_to_spectrum_contravariant}.
\begin{proof}
    Given a map $f \colon R \to S$, we can define $f^\ast \colon \Spec S \to \Spec R$ by $q \mapsto f\inv(q)$. We want to claim that $f\inv(q)$ is actually a prime ideal. Suppose $ab \in f\inv(q)$; equivalently, $f(ab) \in q$; since $f$ is a ring map, this is equivalent to saying that $f(a) \in q$ or $f(b) \in q$, so $a \in f\inv(q)$ or $b \in f\inv(q)$. Thus, $f\inv(q) \in \Spec(R)$.

    We now want to check that $f^\ast$ is a continuous map of topological spaces. Take a closed set $z(I) \subset \Spec(R)$; we want to look at $(f^\ast)\inv (z(I))$. By definition,
    \[ (f^\ast)\inv (z(I)) = \{q \in \Spec S \mid f^\ast(q) \in z(I)\}. \]
    $f^\ast(q)$ is $f\inv(q)$, so $f^\ast(q) \in z(I)$ is equivalent to $I \subset f\inv(q)$, which is the same as $f(I) \subset q$. Thus, the above set is simply $z(f(I))$.

    Finally, we want to check that we actually have a contravariant functor. This is really easy, since the map we have is defined by set maps. First, $(\id_R)^\ast = \id_{\Spec R}$; second, if we have a map $g \colon R \to S$ and $f \colon S \to T$, then $f^\ast \colon \Spec T \to \Spec S$ and $g^\ast \colon \Spec S \to \Spec R$, and we have $(f \circ g)^\ast = g^\ast \circ f^\ast$ as desired. Indeed, for a prime ideal $t \subset T$, we have $(f \circ g)^\ast (t) = (f \circ g)\inv (t) = g\inv \circ f\inv(t) = g^\ast \circ f^\ast(t)$.
\end{proof}

{\color{airforceblue} ``Yay, we've created half the definition for an affine scheme so far.''}

{\color{crimson} \begin{theorem} %\label{thm:pullback}
    Suppose $\Phi \colon R \to S$. Then for $f \in R$, we have $\Phi(f) = f_{\abs{\Spec R}} \circ \Phi^\ast$, where we consider $f \colon \Spec R \to \prod_{p \in \Spec(R)} k(p)$.
    \[ \begin{tikzcd}
        S \arrow[d, dashed, no head] & R \arrow[l, "\Phi(f)"] \arrow[d, dashed, no head] \\
        \Spec S \arrow[r, "\Phi^\ast"] & \Spec R
    \end{tikzcd} \]
\end{theorem}
\begin{proof}
    Hunter said he would leave this as an exercise once he figures out the statement (why...).
\end{proof}}

\begin{theorem} \label{thm:closure_of_set_in_spec}
    Let $Y \subset \Spec R$. Then $\ol Y = Z \left(\bigcap_{p \in Y} p\right)$,\footnote{i swear to god if i should've been using $\kp$ instead of $p$ this whole time i will cry}
\end{theorem}
\begin{proof}
    $\ol Y = z(I)$. Suppose $z(J) \supset Y$; equivalently, $J$ lies in each prime ideal of $Y$, which we can write as
    \[ J \subset \bigcap_{p \in Y} p. \qedhere \]
\end{proof}
\begin{corollary} \label{cor:closure_of_prime_ideal_in_spec}
    $\ol{\{p\}} = z(p) = \{q \in \Spec R \mid p \subset q\}$.
\end{corollary}
Intuitively, ``$p$ is a generic point inside its closure $\ol{\{p\}}$''. As an example from last class, if we look at $\Spec \CC[x, y]$, then $\ol{\{f\}}$ is given by the points on the curve $\bigcup \{(f)\}$. He then said something about the map $\CC[x, y]/(f) \to \Frac(\CC[x, y] / (f))$, but I didn't catch it.

\begin{corollary}
    If $\ol z \subset z(p)$, then every open set containing $\ol z$ contains $p$.
\end{corollary}

As a corollary to our discussion, a prime ideal $p$ is closed if and only if $p$ is maximal.

\newpage
Let's look at an example of an interesting ring that'll we'll see a lot in the future. Recall that if $K$ is a field, then $\Spec K$ is a single point $\{(0)\}$. The next most complicated space, as spectrums go, is a local ring.
\begin{definition} \label{defn:local_ring}
    A \textit{local ring} is a ring with exactly one maximal ideal.
\end{definition}
\begin{corollary} \label{cor:spec_of_local_ring}
    $\Spec R$ has a unique closed point if $R$ is a local ring.
\end{corollary}

As an example, $K[X]_{(x - 5)} = \{\frac{f}{g} : f, g \in K[X], (x - 5) \text{ divides $f$ as many times as it divides $g$}\}$. $(x - 5) \subset K[X]_{(x - 5)}$ is a maximal ideal because the kernel of the evaluation map $\ev_5 \colon K[X]_{(x - 5)} \to K$ is $(x - 5)$. Note that, if $f(5) \neq 0$, then $f$ is a unit. Now,
\[ \frac{f}{g} = (x - 5)^a \frac{\tilde f}{\tilde g}, \]
where $\tilde f$ and $\tilde g$ don't vanish at $5$. Indeed, $\frac{1}{f} \in K[X]_{(x - 5)}$.

Given any ideal $I \subset K[X]_{(x - 5)}$ generated by $\SA$, replace each element of $\SA$ with a power of $x - 5$; then $I$ is generated by some list of powers of $x - 5$. This means $I = (x - 5)^a$, or $I = (a)$, meaning if $a \geq 2$, then it is not prime, since $(x - 5) (x - 5)^{a-1} \in ((x - 5)^a)$. In this manner, in
\[ \Spec K[X]_{(x-5)} = \{(0), (x - 5)\}, \]
the closed sets are $(0)$ (i.e., the whole space), $(x - 5)$, and $\emptyset$.

We then gave a small preview of the structure sheaf. We now end with a final ``easy observation''.

\begin{theorem} \label{thm:embed_with_subspace_topology}
    $\Spec R/I$ embeds into $\Spec R$ with the subspace topology as $z(I)$.
\end{theorem}

``Next week, we're going to do the real meat and potatoes of... spectrums.'